\textbf{Время:} резерв

\textbf{Тезис:} общие пары связывают числитель и знаменатель и требуют учёта структуры событий.

\textbf{Визуал:} среднее независимых весов → ковариация отношения → пересборка событий.

\textbf{Текст:} Если сумма весов делится на число тех же пар, это выборочное среднее. При независимых весах и фиксированном n его дисперсия оценивается выборочной дисперсией весов, делённой на n. В раскрытой записи возникает знаменатель n на n минус один.

Если все веса единичные, средний вес в точности единица. Условная ошибка по весам нулевая — в отличие от результата механической подстановки в формулу двух независимых пуассоновских счётов.

Для общего отношения нужен ковариационный член. Однако его наличие ещё не учитывает повторение треков в разных парах. Конечные неопределённости проверяют на уровне событий с повторением полного анализа. Оценённая из данных нормировка тоже приносит дисперсию и ковариации.
